Dues càrregues elèctriques de $0,03\ \mathrm{\mu C}$ cadascuna, però de signe contrari, es troben separades 40,0~cm.

\begin{apartat}{1}
Representeu i calculeu el vector del camp elèctric en el punt que forma un triangle equilàter amb la posició de les càrregues. Calculeu també el potencial elèctric en el mateix punt.
\end{apartat}

\begin{apartat}{1}
Si modifiquem la distància entre les càrregues fins a duplicar-la, en quant varia l'energia potencial elèctrica de la distribució de càrregues? Expliqueu raonadament si augmenta o disminueix.
\end{apartat}

\blocetiquetat{Dada:}{$k = \dfrac{1}{4\pi\varepsilon_0} = 8,99\times 10^{9}\ \mathrm{N\,m^2\,C^{-2}}$}
